% Resumo em LaTeX: Referenciais, Centro de Massa e Sistemas de Coordenadas
% Compilar com pdflatex (pgfplots necessário)
\documentclass[a4paper,12pt]{article}
\usepackage[T1]{fontenc}
\usepackage[utf8]{inputenc}
\usepackage[portuguese]{babel}
\usepackage{amsmath,amssymb}
\usepackage{siunitx}
\usepackage{graphicx}
\usepackage{caption}
\usepackage{tikz}
\usepackage{tikz-3dplot} 
\usepackage{3dplot}
%\usepackage{pgfplots}
%\pgfplotsset{compat=1.18}
%\usepackage{float}
%\usepackage{hyperref}
\hypersetup{colorlinks=true,linkcolor=blue,urlcolor=blue}
\sisetup{detect-all}

\title{Resumo: Referencial, Centro de Massa, Sistemas de Coordenadas e Trajetórias}
\author{José Gonçalves}
\date{\today}

\begin{document}
	\maketitle
	
	\tableofcontents
	\bigskip
	
	\section{Referencial}
	Um \emph{referencial} ou sistema de referência é um conjunto de eixos e um observador em relação aos quais se descrevem posições, velocidades e acelerações. Em problemas de mecânica, escolher um referencial significa fixar:
	\begin{itemize}
		\item a origem (ponto onde as coordenadas valem zero),
		\item a orientação dos eixos (direções positivas),
		\item a unidade de tempo (relógio) e unidade de comprimento.
	\end{itemize}
	As equações do movimento dependem do referencial — por exemplo, o mesmo movimento pode ter componentes diferentes em referenciais em repouso ou em movimento relativo.
	
	\section{Coordenadas cartesianas}
	\subsection{1D}
	Em \emph{uma dimensão} (reta), a posição de uma partícula é dada por um único número $x(t)$ relativo a uma origem. Exemplos:
	\begin{align*}
		x(t) & = x_0 + v\,t \quad\text{(movimento retilíneo uniforme)} \\
		x(t) & = x_0 + v_0 t + \tfrac{1}{2} a t^2 \quad\text{(movimento com aceleração constante)}
	\end{align*}
	Neste caso, o referencial é unidimensional e representa-se como indicado na fig. \ref{fig:1D}.
	
	\begin{figure}[H]
		\centering
		\begin{tikzpicture}[scale=1]
			% axes
			\draw[->] (0,0) -- (4,0) node[right]{$x$/m};
			% point
			\coordinate (P) at (2,0);
			\fill (P) circle (1.5pt) node[above right]{$P$};
			\coordinate (O) at (0,0);
			\fill (O) circle (1.5pt) node[below]{$O$};
		\end{tikzpicture}
		\caption{Sistema de coordenadas cartesiano (esquemático) a uma dimensão, onde está representado a posição do ponto $P$ num determinado instante.}
		\label{fig:1D}
	\end{figure}
	
	
	
	\subsection{2D}
	No plano, usa-se um par ordenado $(x,y)$. 
	Neste caso, o referencial é bidimensional e representa-se como indicado na fig. \ref{fig:2D}.
	
	\begin{figure}[H]
		\centering
		\begin{tikzpicture}[scale=1]
			% axes
			\draw[->] (0,0) -- (4,0) node[right]{$x$/m};
			\draw[->] (0,0) -- (0,4) node[left]{$y$/m};
			% point
			\coordinate (P) at (2,2);
			\fill (P) circle (1.5pt) node[above right]{$P$};
			\coordinate (O) at (0,0);
			\fill (O) circle (1.5pt) node[below]{$O$};
			\fill (0,2) circle (1pt) node[left]{1,5};
			\fill (2,0) circle (1pt) node[below]{1,5};
			\draw[dashed] (0,2)--(2,2);
			\draw[dashed] (2,0)--(2,2);
			%vetor r
			\draw[->, blue] (0,0) -- (2,2) node[below right]{$\vec{r}$};
		\end{tikzpicture}
		\caption{Sistema de coordenadas cartesiano (esquemático) a duas dimensões, onde está representado a posição do ponto $P$ num determinado instante.}
		\label{fig:2D}
	\end{figure}
	
	O vetor $r$ representa o deslocamento feito desde a oridem até ao ponto $P$ e é dado por: $r = 1,5 \vec{e_x} + 1,5 \vec{e_y}$ (m).
	
	
	\subsection{3D}
	No espaço, as coordenadas são $(x,y,z)$. Neste caso, o referencial é tridimensional e representa-se como indicado na fig. \ref{fig:3D}.
	
	\begin{figure}[H]
		\centering
		\tdplotsetmaincoords{60}{120} 
		\begin{tikzpicture} [scale=3, tdplot_main_coords, axis/.style={->,blue,thick}, 
			vector/.style={-stealth,red,very thick}, 
			vector guide/.style={dashed,red,thick}]
			
			%standard tikz coordinate definition using x, y, z coords
			\coordinate (O) at (0,0,0);
			
			%tikz-3dplot coordinate definition using x, y, z coords
			
			\pgfmathsetmacro{\ax}{0.8}
			\pgfmathsetmacro{\ay}{0.8}
			\pgfmathsetmacro{\az}{0.8}
			
			\coordinate (P) at (\ax,\ay,\az);
			\fill (P) circle (1.5pt) node[above right]{$P$};
			%draw axes
			\draw[axis] (0,0,0) -- (1,0,0) node[anchor=north east]{$x$};
			\draw[axis] (0,0,0) -- (0,1,0) node[anchor=north west]{$y$};
			\draw[axis] (0,0,0) -- (0,0,1) node[anchor=south]{$z$};
			
			%draw a vector from O to P
			%\draw[vector] (O) -- (P);
			
			%draw guide lines to components
			\draw[vector guide]         (O) -- (\ax,\ay,0);
			\draw[vector guide] (\ax,\ay,0) -- (P);
			\draw[vector guide]         (P) -- (0,0,\az);
			\draw[vector guide] (\ax,\ay,0) -- (0,\ay,0);
			\draw[vector guide] (\ax,\ay,0) -- (0,\ay,0);
			\draw[vector guide] (\ax,\ay,0) -- (\ax,0,0);
			\node[tdplot_main_coords,anchor=east]
			at (\ax,0,0){(\ax, 0, 0)};
			\node[tdplot_main_coords,anchor=west]
			at (0,\ay,0){(0, \ay, 0)};
			\node[tdplot_main_coords,anchor=south]
			at (0,0,\az){(0, 0, \az)};
			
			%vector
			\draw[->, orange] (0,0,0) -- (\ax,\ay,\az) node[above left]{$\vec{r}$};
		\end{tikzpicture}
		\caption{Sistema de coordenadas cartesiano (esquemático) a três dimensões, onde está representado a posição do ponto $P$ num determinado instante.}
		\label{fig:3D}
	\end{figure}
	
	O vetor $r$ representa o deslocamento feito desde a oridem até ao ponto $P$ e é dado por: $r = 0,8 \vec{e_x} + 0,8 \vec{e_y} + 0,8 \vec{e_z}$ (m).
	
	
	
	\section{Coordenadas cilíndricas}
	Coordenadas cilíndricas são úteis quando existe simetria em torno de um eixo (por exemplo, problemas envolvendo cilindros ou eixos).
	
	%Cylindrical Coordinate
	\newcommand{\cylindricalcoordinate}[4]{%
		\coordinate (#4) at ({#1*cos(#2)},{#1*sin(#2)},{#3});%
		\coordinate (#4xy) at ({#1*cos(#2)},{#1*sin(#2)},0);%
	}
	%Axis Angles
	\tdplotsetmaincoords{70}{110}
	
	%Macros
	\pgfmathsetmacro{\rvec}{5}
	\pgfmathsetmacro{\phivec}{45}
	\pgfmathsetmacro{\zvec}{4}
	
	\pgfmathsetmacro{\drvec}{1.5}
	\pgfmathsetmacro{\dphivec}{20}
	\pgfmathsetmacro{\dzvec}{1}
	
	%Layers
	\pgfdeclarelayer{background}
	\pgfdeclarelayer{foreground}
	
	\pgfsetlayers{background, main, foreground}
	
	\begin{tikzpicture}[tdplot_main_coords]
		\centering
		%Coordinates
		\coordinate (O) at (0,0,0);
		%%
		\cylindricalcoordinate{\rvec}{\phivec}{\zvec}{A}
		\cylindricalcoordinate{\rvec}{\phivec}{(\zvec+\dzvec)}{B}
		\cylindricalcoordinate{\rvec}{(\phivec+\dphivec)}{\zvec}{C}
		\cylindricalcoordinate{\rvec}{(\phivec+\dphivec)}{(\zvec+\dzvec)}{D}
		%%
		\cylindricalcoordinate{(\rvec+\drvec)}{\phivec}{\zvec}{A'}
		\cylindricalcoordinate{(\rvec+\drvec)}{\phivec}{(\zvec+\dzvec)}{B'}
		\cylindricalcoordinate{(\rvec+\drvec)}{(\phivec+\dphivec)}{\zvec}{C'}
		\cylindricalcoordinate{(\rvec+\drvec)}{(\phivec+\dphivec)}{(\zvec+\dzvec)}{D'} 
		
		%Axis
		\begin{pgfonlayer}{background}
			\draw[thick,-latex] (0,0,0) -- (6,0,0) node[pos=1.1]{$x$};
			\draw[thick,-latex] (0,0,0) -- (0,6,0) node[pos=1.05]{$y$};
			\draw[thick,-latex] (0,0,0) -- (0,0,7) node[pos=1.05]{$z$};
		\end{pgfonlayer}
		
		%Vectors
		\begin{pgfonlayer}{main}
			\draw[blue, thick] (O) -- (A);
			\draw[thick] (O) -- (Axy) node [pos=0.6, below left] {$\varpi$};
			\draw (A) -- ($(A)-(Axy)$) node [left] {$z$};
			\draw (B) -- ($(A)-(Axy)+(0,0,\dzvec)$) node [left] {$z+\mathrm{d}z$};
			\draw (D) -- ($(A)-(Axy)+(0,0,\dzvec)$);
		\end{pgfonlayer}
		\begin{pgfonlayer}{background}
			\draw[dashed] (C) -- ($(A)-(Axy)$) node [left] {$z$};
		\end{pgfonlayer}
		
		%Help Lines
		\begin{pgfonlayer}{background}
			\draw (A) -- (Axy);
			\draw (A') -- (A'xy);
			\draw[thick] (Axy) -- (A'xy) node [pos=0.6, below left] {$\mathrm{d}\varpi$};
			%
			\draw (O) -- (D'xy);
			\draw[dashed] (C) -- (Dxy); 
			\draw (C') -- (C'xy);
			%%Arcs
			\tdplotdrawarc{(0,0,0)}{\rvec}{\phivec}{\phivec+\dphivec}{}{}
			\tdplotdrawarc{(0,0,0)}{\rvec+\drvec}{\phivec}{\phivec+\dphivec}{}{}
		\end{pgfonlayer}
		
		%Angles
		\begin{pgfonlayer}{foreground}
			%Phi, dPhi
			\tdplotdrawarc[-stealth]{(O)}{0.9}{0}{\phivec}{anchor=north}{$\phi$}
			\tdplotdrawarc[-stealth]{(O)}{1.5}{\phivec}{\phivec + \dphivec}{}{}
			\node at (1.4,1.9,0) {$\mathrm{d}\phi$};    
		\end{pgfonlayer}
		
		%Differential Volume
		
		%%Lines
		\begin{pgfonlayer}{foreground}
			\draw[thick] (A) -- (B) -- (B') -- (A') -- cycle node [midway, below] {$\mathrm{d}\varpi$};
			\draw[thick] (D) -- (D') -- (C') node [midway, right] {$\mathrm{d}z$};
		\end{pgfonlayer}
		\begin{pgfonlayer}{background}
			\draw[thick, dashed] (C') -- (C) -- (D);
		\end{pgfonlayer}
		
		%%Curved
		\begin{pgfonlayer}{background}
			\tdplotdrawarc[thick, dashed]{(0,0,\zvec)}{\rvec} {\phivec}{\phivec+\dphivec}{}{}
		\end{pgfonlayer}
		\begin{pgfonlayer}{foreground}
			\tdplotdrawarc[thick]{(0,0,\zvec+\dzvec)}{\rvec} {\phivec}{\phivec+\dphivec}{}{}
			\tdplotdrawarc[thick]{(0,0,\zvec+\dzvec)}{\rvec+\drvec} {\phivec}{\phivec+\dphivec}{}{}
			\tdplotdrawarc[thick]{(0,0,\zvec)}{\rvec+\drvec} {\phivec}{\phivec+\dphivec}{below}{$\varpi\mathrm{d}\phi$}
		\end{pgfonlayer}
		
		%%Fill Color
		\begin{pgfonlayer}{main}
			%Front
			\fill[black, opacity=0.15] (B) to (B') to[bend right=4] (D') to (D) to[bend left=4] cycle;
			\fill[black, opacity=0.6] (B') to[bend right=4] (D') to (C') to[bend left=4] (A') to cycle;
			\fill[black, opacity=0.4] (B) to (B')  to (A') to (A) to cycle;
		\end{pgfonlayer}
		\begin{pgfonlayer}{background}
			%Back
			\fill[black!50, opacity=0.5] (D) to (D') to (C') to (C) to cycle;
			\fill[black!50, opacity=0.5] (B) to[bend right=4] (D) to (C) to[bend left=4] (A) to cycle;
			\fill[black!50, opacity=0.5] (A) to (A') to[bend right=4] (C') to (C) to[bend left=4] cycle;
		\end{pgfonlayer}
		
	\end{tikzpicture}
	
	\section{Coordenadas esféricas}
	Coordenadas esféricas são usadas quando existe simetria esférica.
	
	%Angle Definitions
	%-----------------
	
	%set the plot display orientation
	%synatax: \tdplotsetdisplay{\theta_d}{\phi_d}
	\tdplotsetmaincoords{60}{110}
	
	%define polar coordinates for some vector
	%TODO: look into using 3d spherical coordinate system
	\pgfmathsetmacro{\rvec}{.8}
	\pgfmathsetmacro{\thetavec}{45}
	\pgfmathsetmacro{\phivec}{60}
	
	%start tikz picture, and use the tdplot_main_coords style to implement the display 
	%coordinate transformation provided by 3dplot
	\begin{tikzpicture}[scale=5,tdplot_main_coords]
		
		\shadedraw[tdplot_screen_coords,ball color = white] (0,0) circle (\rvec);
		
		%set up some coordinates 
		%-----------------------
		\coordinate (O) at (0,0,0);
		
		%determine a coordinate (P) using (r,\theta,\phi) coordinates.  This command
		%also determines (Pxy), (Pxz), and (Pyz): the xy-, xz-, and yz-projections
		%of the point (P).
		%syntax: \tdplotsetcoord{Coordinate name without parentheses}{r}{\theta}{\phi}
		\tdplotsetcoord{P}{\rvec}{\thetavec}{\phivec}
		
		%draw figure contents
		%--------------------
		
		%draw the main coordinate system axes
		\draw[thick,->] (0,0,0) -- (1,0,0) node[anchor=north east]{$x''$};
		\draw[thick,->] (0,0,0) -- (0,1,0) node[anchor=north west]{$y''$};
		\draw[thick,->] (0,0,-1) -- (0,0,1) node[anchor=south]{$z''$};
		
		%draw a vector from origin to point (P) 
		\draw[-stealth,color=black] (O) -- (P) node[midway,above] {$r$};
		
		%draw projection on xy plane, and a connecting line
		\draw[dashed, color=red] (O) -- (Pxy);
		\draw[dashed, color=red] (P) -- (Pxy);
		
		%draw the angle \phi, and label it
		%syntax: \tdplotdrawarc[coordinate frame, draw options]{center point}{r}{angle}{label options}{label}
		\tdplotdrawarc{(O)}{0.2}{0}{\phivec}{anchor=north}{$\alpha$}
		
		
		%set the rotated coordinate system so the x'-y' plane lies within the
		%"theta plane" of the main coordinate system
		%syntax: \tdplotsetthetaplanecoords{\phi}
		\tdplotsetthetaplanecoords{\phivec}
		
		%draw theta arc and label, using rotated coordinate system
		\tdplotdrawarc[tdplot_rotated_coords]{(0,0,0)}{0.5}{\thetavec}{90}{anchor=south west}{$\beta$}
		
		%draw some dashed arcs, demonstrating direct arc drawing
		\draw[thin,tdplot_rotated_coords] (\rvec,0,0) arc (0:180:\rvec);
		
		
		
	\end{tikzpicture}
	
	\section{Trajetória}
	\textbf{Definição:} a trajetória de uma partícula é o conjunto de pontos no espaço ocupados por ela (posições) ao longo do tempo — geometricamente, é a curva \footnote{Uma curva em termos físicos, também pode ser representada por uma reta se o vetor posição não mudar de direção.} descrita pelo vetor posição $\mathbf{r}(t)$. Em termos paramétricos:
	\[ \text{trajetória} = \{\,\mathbf{r}(t) : t_1\le t \le t_2\,\}. \]
	Uma trajetória pode ser descrita em coordenadas cartesianas $(x(t),y(t),z(t))$ ou em coordenadas adequadas ao problema (cilíndricas, esféricas, etc.).
	Neste capítulo, vamos apenas abordar movimentos unidimensionais. Assim, o vetor posição $\mathbf{r}(t)$ será apenas descrito apenas por uma coordenada, $x(t)$.
	
	
	VOU AQUI
	
	
	
	\section{Centro de massa}
	\subsection{Definição geral}
	Para um sistema de $N$ partículas com massas $m_i$ e posições $\mathbf{r}_i$, o centro de massa (CM) é:
	\begin{equation}\label{eq:CM_discrete}
		\mathbf{R}_{\mathrm{CM}} = \frac{1}{M} \sum_{i=1}^N m_i\,\mathbf{r}_i,\qquad M=\sum_{i=1}^N m_i.
	\end{equation}
	Para um corpo contínuo com densidade $\rho(\mathbf{r})$, usamos a integral:
	\begin{equation}\label{eq:CM_cont}
		\mathbf{R}_{\mathrm{CM}} = \frac{1}{M}\int_{V} \mathbf{r}\,\rho(\mathbf{r})\,dV,\qquad M=\int_{V}\rho(\mathbf{r})\,dV.
	\end{equation}
	
	\subsection{Exemplo simples}
	Duas partículas: massa $m_1$ em $x_1$ e massa $m_2$ em $x_2$ (1D). Então
	\[ X_{\mathrm{CM}} = \frac{m_1 x_1 + m_2 x_2}{m_1+m_2}. \]
	\textbf{Exemplo numérico:} $m_1=2\,\mathrm{kg}$ em $x_1=1\,\mathrm{m}$ e $m_2=3\,\mathrm{kg}$ em $x_2=5\,\mathrm{m}$. Então
	\begin{align*}
		X_{\mathrm{CM}} &= \frac{2\cdot 1 + 3\cdot 5}{2+3} = \frac{2+15}{5} = \frac{17}{5} = 3.4\,\mathrm{m}.
	\end{align*}
	
	\subsection{Propriedades importantes}
	\begin{itemize}
		\item O centro de massa se comporta como se toda a massa estivesse concentrada nele quando consideramos forças externas (segunda lei de Newton aplicada ao CM): $M\mathbf{A}_{\mathrm{CM}}=\sum \mathbf{F}_{\mathrm{ext}}$.
		\item Para colisões e problemas de impulso, é conveniente analisar o movimento do CM.
	\end{itemize}
	
	\section{Trajetórias: retilíneas e curvilíneas}
	\subsection{Trajetória retilínea}
	Uma trajetória retilínea é aquela em que o movimento ocorre ao longo de uma reta fixa. Em 1D é sempre retilíneo; em 2D/3D acontece quando $\mathbf{r}(t)$ permanece colinear a uma direção fixa.
	Exemplo: movimento uniformemente acelerado em linha reta:
	\[ x(t)=x_0+v_0 t + \tfrac{1}{2} a t^2. \]
	
	\subsection{Trajetória curvilínea}
	Quando a trajetória é uma curva genérica (circunferência, parábola, hélice, etc.). Em coordenadas paramétricas no plano:
	\[ x(t)=f(t),\quad y(t)=g(t). \]
	A curvatura e o raio de curvatura são grandezas que descrevem a ``curvatura'' local da trajetória.
	
	\paragraph{Exemplos clássicos}
	\begin{itemize}
		\item Movimento circular uniforme: $x(t)=R\cos(\omega t)$, $y(t)=R\sin(\omega t)$ (trajetória = circunferência).
		\item Projéteis (sem resistência do ar): trajetória parabólica no plano $x$--$y$.
		\item Hélice: $x(t)=R\cos(\omega t)$, $y(t)=R\sin(\omega t)$, $z(t)=v_z t$.
	\end{itemize}
	
	\section{Gráficos posição---tempo (exemplos)}
	Gráficos posição---tempo mostram $x$ (ou $y$, ou componente de posição) em função do tempo. Vamos apresentar dois exemplos com códigos \texttt{pgfplots}:
	\begin{figure}[H]
		\centering
		\begin{tikzpicture}
			\begin{axis}[xlabel={$t$ (s)},ylabel={$x$ (m)},title={Movimento retilíneo uniforme: $x(t)=x_0+v t$},width=10cm,height=6cm]
				\addplot[domain=0:10,samples=2]{2 + 1.5*x};
				\addlegendentry{$x(t)=2+1.5t$}
			\end{axis}
		\end{tikzpicture}
		\caption{Exemplo: posição linear no tempo (velocidade constante 1.5 m/s).}
	\end{figure}
	
	\begin{figure}[H]
		\centering
		\begin{tikzpicture}
			\begin{axis}[xlabel={$t$ (s)},ylabel={$x$ (m)},title={Movimento com acelera\c{c}\~ao constante: $x(t)=x_0+v_0 t + \tfrac{1}{2}a t^2$},width=10cm,height=6cm]
				\addplot[domain=0:10,samples=201]{1 + 2*x + 0.5*0.8*x^2};
				\addlegendentry{$x(t)=1+2t+0.4t^2$}
			\end{axis}
		\end{tikzpicture}
		\caption{Exemplo: movimento com acelera\c{c}\~ao constante (a=0.8 m/s$^2$).}
	\end{figure}
	
	\section{Notas sobre interpretação dos gráficos}
	\begin{itemize}
		\item A inclinação (derivada) do gráfico posição---tempo dá a velocidade instantânea: $v(t)=dx/dt$. Para o gráfico linear a inclinação é constante (velocidade constante).
		\item A concavidade do gráfico posição---tempo indica aceleração: se o gráfico é convexo para cima (concavidade para cima), a aceleração é positiva.
		\item Um ponto com inclinação zero corresponde a velocidade instantânea nula (máximo ou mínimo local da posição).
	\end{itemize}
	
	\section{Apêndice: pequenas demonstra\c{c}\~oes e fórmulas úteis}
	\subsection{Velocidade média e instantânea}
	Velocidade média entre $t_1$ e $t_2$:
	\[ v_{\mathrm{med}}=\frac{x(t_2)-x(t_1)}{t_2-t_1}. \]
	Velocidade instantânea: derivada temporal
	\[ v(t)=\lim_{\Delta t\to 0}\frac{x(t+\Delta t)-x(t)}{\Delta t}=\frac{dx}{dt}. \]
	
	\subsection{Exercício resolvido (rápido)}
	Uma partícula tem $x(t)=3t^2-2t+1$ (m). Determine $v(t)$ e $a(t)$.
	\begin{align*}
		v(t) &= \frac{dx}{dt} = 6t-2 \quad(\mathrm{m/s}),\\
		a(t) &= \frac{d^2x}{dt^2} = 6 \quad(\mathrm{m/s^2}).
	\end{align*}
	
	\bigskip
%	\noindent Se quiser, posso:
%	\begin{itemize}
%		\item adicionar mais exemplos (projétil, movimento circular),
%		\item incluir figuras 3D ou animações em beamer (PDF),
%		\item exportar o ficheiro como PDF (se preferires, pede que eu gere o PDF). 
%	\end{itemize}
	
\end{document}
